\section{Las aventuras salvajes de un joven íntegro: de la App Store al NLP}

La sección anterior trató lo particular de la última entrevista; esta sección vuelve al punto de partida de Peak. Su trayectoria temprana se parece a una historia de la tecnología en miniatura. En la preparatoria desarrollaba un navegador de terceros para iOS, aprovechó la bonanza de los primeros años de la App Store y obtuvo flujo de efectivo con el modelo más simple, el buy-copy. Esa experiencia le hizo entender muy pronto que un cambio de plataforma puede poner por un tiempo a los desarrolladores independientes, a las empresas emergentes y a las grandes compañías en la misma línea de salida.

\begin{figure}[H]
\centering
\includegraphics[width=0.96\textwidth]{manus-journey.png}
\caption{La travesía fantástica de Manus: de la App Store, el NLP y Open IE hasta Manus y la adquisición por parte de Meta.}
\end{figure}

\begin{knowledgebox}{Lectura de la figura: Manus no surgió de la nada}
Cada nodo de la figura nace de una necesidad concreta. El navegador tenía que ser más rápido con una red débil, de ahí el estudio de cómo predecir el siguiente clic; la predicción de clics llevó al NLP; el NLP llevó a Word2Vec, a los grafos de conocimiento y a Open IE; toda esa experiencia desembocó finalmente en la conversión en producto de los Agent.
\end{knowledgebox}

\subsection{La etapa salvaje de la App Store}

Esta subsección examina la primera ventana de plataforma. Peak cuenta que la App Store le dio un indicador externo: el software que armaba experimentando por su cuenta podía generar valor económico. Ese indicador era muy importante para un estudiante de preparatoria, porque convertía lo que parecía «descuidar lo que debía hacer» en una capacidad de producto verificable. Los nuevos medios, los nuevos canales de distribución y los nuevos modelos de negocio de los inicios del internet móvil dejaron una ventana abierta para los desarrolladores independientes.

\begin{figure}[H]
\centering
\includegraphics[width=0.94\textwidth]{appstore-to-ai-platform.png}
\caption{La etapa salvaje de la App Store frente a la era de la IA: diferencias en la bonanza de plataforma.}
\end{figure}

\begin{knowledgebox}{Lectura de la figura: un avance técnico en IA no equivale a una nueva bonanza de plataforma}
A la izquierda, la etapa salvaje de la App Store trajo un cambio en el medio de hardware y una bonanza de distribución; a la derecha, la era de la IA presenta avances técnicos enormes, pero no una nueva plataforma igual de clara, de modo que los gigantes, las empresas emergentes y los desarrolladores independientes reaccionan con la misma rapidez y la ventana es más estrecha.
\end{knowledgebox}

\subsection{Del precargado del navegador al NLP}

En la época del navegador, quería predecir el siguiente clic del usuario y cargar por adelantado la página siguiente, para mejorar la experiencia con una red débil. Esa necesidad concreta lo llevó al NLP. Más tarde, Word2Vec le mostró que el lenguaje natural podía convertirse en vectores densos, y los grafos de conocimiento lo llevaron a intentar responder las preguntas de los usuarios de forma estructurada.

\begin{figure}[H]
\centering
\includegraphics[width=0.96\textwidth]{nlp-lineage.png}
\caption{Del precargado del navegador al NLP: una necesidad concreta orienta la ruta de aprendizaje.}
\end{figure}

\begin{knowledgebox}{Lectura de la figura: la ruta de aprendizaje la orientan los problemas}
En la figura, el paso del navegador al NLP no consiste en elegir primero un sector y luego buscar un problema, sino en partir de una necesidad concreta: «¿cómo ser más rápido con una red débil?». El buen aprendizaje técnico rara vez consiste en perseguir modas en abstracto; suele estar orientado por problemas reales de producto.
\end{knowledgebox}

\subsection{Por qué el navegador no conviene a un emprendedor a largo plazo}

Peak también ha reflexionado que los navegadores de terceros nunca han sido un buen proyecto de largo plazo para un emprendedor. El navegador se parece más a un producto de entrada que hace un gigante cuando ya tiene canales de distribución, porque por naturaleza exige distribución a nivel de sistema, entrada predeterminada, colaboración con el ecosistema y confianza en la seguridad. Este juicio también sirve para pensar en los navegadores con IA de hoy: un navegador con IA parece una oportunidad para emprender, pero sin distribución ni una entrada en el sistema, una empresa emergente queda fácilmente aplastada por los gigantes.

La relación entre Manus y el navegador es más sutil. Manus no se limita a hacer un navegador, sino que busca que el Agent actúe dentro del navegador y del mundo de las páginas web. Es decir, no necesita convertirse en el navegador mismo, sino en la capa de acción que atraviesa páginas web, herramientas y tareas. Esto explica por qué hereda la preocupación por los problemas que Peak tuvo en su etapa temprana de navegador/NLP y, a la vez, no regresó al camino de «hacer un navegador».

\begin{longtable}{p{0.24\textwidth}p{0.34\textwidth}p{0.33\textwidth}}
\toprule
Objeto & Tarea central & Dificultad para emprender \\
\midrule
Navegador tradicional & Mostrar páginas web, administrar pestañas, alojar la búsqueda y la entrada a los sitios web & La entrada predeterminada y el control de la distribución suelen estar en manos de los gigantes. \\
Navegador con IA & Añadir al navegador funciones como resumir, responder preguntas y llenar formularios automáticamente & Si solo mejora el navegador, sigue sujeto a las restricciones de entrada y de ecosistema. \\
Capa de acción del Agent & Completar tareas a través de páginas web, herramientas, archivos y cuentas & Requiere permisos, pagos, colaboración frente a los mecanismos contra el rastreo automatizado y ejecución confiable. \\
\bottomrule
\end{longtable}

\begin{importantbox}{De la entrada del navegador a la capa de acción del Agent}
El navegador se encarga de mostrar las páginas web; el Agent, de entender el objetivo y actuar. Si un navegador con IA es solo «un navegador más inteligente», puede seguir limitado por la entrada; si un Agent puede ejecutar tareas a través de herramientas, puede convertirse en una capa de acción por encima del navegador.
\end{importantbox}

\subsection{Resumen de la sección}

El carácter de fondo de Manus proviene de la experiencia temprana como desarrollador independiente: primero resolver un problema concreto y luego aprender la tecnología siguiendo ese problema. Esta trayectoria explica por qué Peak da más peso a la retroalimentación real del producto que a los grandes relatos sobre la IA.
